\chapter{Unit 1: Exploring One-Variable Data}

\section{Essential Concepts \& Cheat Sheet}

When describing any quantitative single-variable distribution on the AP Exam, you \textbf{must} address four key characteristics: \textbf{SOCS} (Shape, Outliers, Center, Spread) always written \textbf{in context} of the problem.

\begin{formulabox}[Core Formulas: One-Variable Statistics]
\textbf{1. Mean ($\bar{x}$):}
\[
\bar{x} = \frac{\sum x_i}{n} = \frac{1}{n} \sum x_i
\]

\textbf{2. Sample Standard Deviation ($s_x$):}
\[
s_x = \sqrt{\frac{1}{n-1}\sum (x_i - \bar{x})^2}
\]

\textbf{3. Standardized Score ($z$-score):}
\[
z = \frac{x - \mu}{\sigma} \quad \text{or} \quad z = \frac{x - \bar{x}}{s_x}
\]
Represents how many standard deviations a value $x$ falls above ($z > 0$) or below ($z < 0$) the mean.
\end{formulabox}

\section{Outlier Identification Rules}

\subsection{The 1.5 $\times$ IQR Rule (for Skewed Data / Medians)}
\begin{itemize}
    \item Calculate Interquartile Range: $\text{IQR} = Q_3 - Q_1$
    \item Lower Outlier Boundary: $\text{Lower Fence} = Q_1 - 1.5(\text{IQR})$
    \item Upper Outlier Boundary: $\text{Upper Fence} = Q_3 + 1.5(\text{IQR})$
    \item Any data point $x < \text{Lower Fence}$ or $x > \text{Upper Fence}$ is an \textbf{outlier}.
\end{itemize}

\subsection{The 2-Standard-Deviation Rule (for Symmetric / Normal Data)}
\begin{itemize}
    \item Outliers are points beyond: $\mu \pm 2\sigma$ (or $\mu \pm 3\sigma$ for extreme outliers).
\end{itemize}

\section{Resistant vs. Non-Resistant Summary Statistics}

\begin{table}[H]
\centering
\small
\begin{tabular}{@{}lll@{}}
\toprule
\textbf{Measure} & \textbf{Resistant to Outliers / Skew?} & \textbf{When to Use on Exam} \\
\midrule
\textbf{Median} & \textbf{Yes} (Unaffected by extremes) & Skewed distributions \\
\textbf{IQR} & \textbf{Yes} (Middle 50\% only) & Skewed distributions \\
\textbf{Mean} & \textbf{No} (Pulled toward the tail/outliers) & Symmetric, unimodal distributions \\
\textbf{Standard Deviation} & \textbf{No} (Squared deviations inflate) & Symmetric, unimodal distributions \\
\textbf{Range} & \textbf{No} (Directly depends on Max/Min) & Never use as sole measure of spread \\
\bottomrule
\end{tabular}
\caption{Resistant vs Non-Resistant Measures}
\end{table}

\begin{trapbox}[AP Exam Trap Alert: Comparing Distributions]
When comparing two distributions on an FRQ:
\begin{enumerate}
    \item You \textbf{must} use comparative words (e.g., \emph{"Distribution A has a greater median than Distribution B"}, NOT just \emph{"Distribution A has a median of 5 and B has 3"}).
    \item You \textbf{must} include context (e.g., grams of sugar, test scores, heights in inches).
    \item Never mix resistant and non-resistant measures (compare median to median, or mean to mean).
\end{enumerate}
\end{trapbox}

\section{The Empirical Rule (68-95-99.7 Rule)}
For any approximately Normal distribution $N(\mu, \sigma)$:
\begin{itemize}
    \item $\approx 68\%$ of data falls within $\mu \pm 1\sigma$.
    \item $\approx 95\%$ of data falls within $\mu \pm 2\sigma$.
    \item $\approx 99.7\%$ of data falls within $\mu \pm 3\sigma$.
\end{itemize}
